\section{Summary and Discussion} 
\label{conclusion}

We have collected a diverse panel data set covering 30,000+ popular playlists on Spotify to estimate the causal treatment effect of featuring playlists on Spotify's Search Page, and compare the impact of being featured on this high-traffic page with the effect of adding tracks by popular major label superstar artists to playlists. We find that users are responsive to the featuring of playlists, and our estimates suggest that inclusion on Spotify's Search Page raises daily followers by 0.95\%. For large playlists like ``Reggaeton Classic" (with 1.9 million followers), this implies an increase of 18,050 additional daily followers after being featured in Spotify's app. By comparison, we find the effect of including tracks by popular major label superstars in a playlist to be $0.45\%$. Our analysis reveals that the average treatment effect of a playlist being featured is about two times greater than the effect of adding a track by a ``superstar” major label artist to the list. %Our findings are \textit{robust} as the differences in treatment effects are recovered in different model specifications that address the empirical identification challenges, and \textit{generalizable} as the estimation data includes 30,000+ playlists in 66 music categories.

Our work has important implications for policymakers and music industry professionals.
First, as major music labels have watched their market share on music streaming services erode in recent years, some music executives have started to wonder about the value of their ``superstar" artists.\footnote{See \href{https://www.musicbusinessworldwide.com/as-major-label-market-share-falls-on-spotify-can-we-expect-even-more-indie-investment-from-the-big-three/}{https://www.musicbusinessworldwide.com/as-major-label-market-share-falls-on-spotify-can-we-expect-even-more-indie-investment-from-the-big-three/}, accessed October 31, 2023.} %Some industry insiders have even announced the ``death of the streaming megahit," and major labels have increased investments in alternative distribution services to make up for the loss of market share.\footnote{See \href{https://www.musicbusinessworldwide.com/are-we-witnessing-the-death-of-the-streaming-megahit/}{https://www.musicbusinessworldwide.com/are-we-witnessing-the-death-of-the-streaming-megahit/}, accessed October 31, 2023.} 
However, our research suggests that ``superstars" continue to create value on music streaming services: consumers respond to adding major label superstars with an estimated increase in daily followers of $0.45\%$ in the fifteen days after the event. 
This finding is of interest to record labels who are disputing royalty rates with streaming platforms in the US.\footnote{See case docket \textit{19-CRB-0005-WR} at the Copyright Royalty Board, available at \href{https://app.crb.gov/case/detail/19-CRB-0005-WR\%20\%282021-2025\%29}{https://app.crb.gov/}, accessed March 5, 2024.} Our findings can also aid court judges in assessing the consumer value of major label content on music streaming platforms like Spotify.

Second, we measure the short-term treatment effect of including playlists on the Search Page, a gateway to the app’s content where a select number of playlists --- typically Spotify-curated playlists --- are brought to the user's attention. Our finding that followers increase by $0.95\%$ implies that Spotify is effective in guiding user attention by selectively highlighting playlists on the platform \citep[e.g., ][]{goldfarb_tucker_2019}. Importantly, the fact that the featuring effect is about twice the size of the superstar effect suggests the success of Spotify playlists is much less dependent on the availability of artist-specific (superstar) tracks than it is to being prominently featured on the platform. We believe this difference in effect may in the future drive record labels to pay Spotify to be prominently featured.\footnote{Pay-for-play arrangements are already being used by Spotify in other contexts, such as the recently introduced (``Pay For Influence") Discovery Mode that allows artists and record labels to influence the algorithm that selects what music is played when users use the streaming platform's auto-play or personalized radio function. By opting in, record labels agree to be paid a lower royalty rate for streams. This recent development aligns with our results and showcases the dependency of artists (and record labels more broadly) on Spotify's programming decisions. See \href{https://www.isrc.com/3-Things-To-Know-About-Spotifys-Discovery-Mode.php}{https://www.isrc.com/3-Things-To-Know-About-Spotifys-Discovery-Mode.php}, accessed May 1, 2023.}  

To foster competition among playlist curators, policymakers could legally impose a minimum fraction of non-Spotify playlists that need to be featured on the platform's Search Page on a given day. 
Spotify already features its own curated playlists in almost 90\% of cases (Table \ref{tab:promoplaylist}). A policy requiring Spotify to feature a minimum percentage of non-Spotify-curated playlists would be analogous to recently implemented EU legislation requiring streaming services like Apple TV+ and Netflix to feature at least $30\%$ European content or face a European Union customer block.\footnote{See \href{https://www.gov.ie/en/publication/d8e4c-online-safety-and-media-regulation-bill/}{Online Safety and Media Regulation Bill}, accessed July 5, 2022.}

Our work opens several avenues for inquiry about the economics of streaming. First, our data does not incorporate information on all potential factors influencing users' demand for playlists, such as ``announcements" to advertise playlists or concerts on Spotify. Second, we study playlist followers but lack sufficient data on a playlist's number of listeners. Third, the results extend to the set of 30,000+ most popular playlists on the platform. However, a larger sample of playlists would allow for a better understanding of how Spotify can exert power over music labels through algorithmic design. %While such data is unavailable to us, music labels may have access to suitable, user-level data. Such data would also be ideally suited to model the consecutive consumption of multiple playlists, or differentiating between Spotify's consumption channels (e.g., playlists discovered via the Search Page, vs. playlists stored in users' libraries). 
Fourth, our data primarily allows us to identify short-term treatment effects.
Yet, for music labels, understanding the long-term dynamics of a playlist's follower evolution is important to develop successful playlist brands. For policymakers, in turn, long-term estimates would enable a more nuanced understanding of the power dynamics between digital platforms and their suppliers to devise effective regulatory policies.

%For example, the insignificant short-run estimate of a major label's share of content on playlist followers should not be construed to suggest that artists other than those classified as superstars lack long-term importance in enhancing playlist appeal.


